定义 $\mathcal Lw=M:H(w)=\operatorname{Re}\operatorname{tr}(MH(w))$。由各槽的公式得 $\mathcal Lw=M^\mathbb R:D^2w$。列散度条件给出 $\mathcal Lw=\operatorname{div}(M^\mathbb R\nabla w)$。由于 $\eta$ 只依赖于 $z$，有 $\mathcal L\eta=M_{TT}\Delta_z\eta$。Green 恒等式给出
\[
\int_\Omega u\mathcal L\eta-\eta\mathcal Lu
=-\int_{\partial_{\rm lat}\Omega}\eta(\nu^*M\nu)\partial_{\mathbf n}u\,dS,
\quad
\mathbf n=\nabla r/|\nabla r|,\quad \nu=G(r)/|\nabla r|.
\]
侧壁上含 $u$ 的项为零，而 $\eta$ 的紧支撑使参数域端部的各项消失。由于迹为零，侧壁上有 $\nabla u=(\partial_{\mathbf n}u)\mathbf n$。作光滑除法得侧壁附近的 $u=r\rho$，由此得到 $\partial_{\mathbf n}u=\rho|\nabla r|$。在全空间上应用余面积公式，将 $dS/|\nabla r|$ 改写为 $\delta(r)\,dz\,d\zeta$；固定 $z$ 应用余面积公式，将 $\delta(r)\,d\zeta$ 改写为 $dS_{\rm fib}/|\nabla_\zeta r|$。代入即得所述公式。
